\section{Errors and Exceptions}\label{sec:ErrorsAndExceptions}
Although error handling is more than exception\index{exception} handling, as stated already above, Exceptions\index{exception} are still a topic that needs to be discussed in this context.

This topic has various dimensions:
\begin{itemize}
	\item{The selection of the Exception\index{exception} to throw to report an error condition}
	\item{The definition of custom Exceptions\index{exception}}
	\item{Generating and throwing Exceptions\index{exception}}
	\item{Catching Exceptions\index{exception}}
\end{itemize}

\subsection{Which Exception to throw?}
When your code needs to issue an Exception\index{exception} to signal an Error or a Failure, you must choose the respective Exception\index{exception} class\index{class!exception class}. That selection gives already some basic information about the kind and origin of the problem, through the name of the Exception\index{exception} class\index{class!exception class}, but also through the package that class belongs to. This means that the classes

\begin{itemize}[nosep]
	\item{\lstinline|java.lang.Throwable|\index{java.lang!Throwable} \autocite{ORACLE_DOC_THROWABLE_CLASS}}
	\item{\lstinline|java.lang.Exception|\index{java.lang!Exception} \autocite{ORACLE_DOC_EXCEPTION_CLASS}}
	\item{\lstinline|java.lang.RuntimeException|\index{java.lang!RuntimeException} \autocite{ORACLE_DOC_RUNTIMEEXCEPTION_CLASS}}
	\item{\lstinline|java.lang.Error|\index{java.lang!Error} \autocite{ORACLE_DOC_ERROR_CLASS}}
\end{itemize}

are forbidden, because they are too generic (and they are causing issues when catching\index{try!catch} exceptions\index{exception} what we will discuss later).

When scanning the Javadoc for the Java Runtime Library \autocite{ORACLE_DOC_OVERVIEW}, you may find several Exceptions\index{exception} with names that may fit for your purpose, but when you look after their package, they belong to an area that is by no means related to your method\index{method}.

Currently, our \lstinline|textToNumber()| method\index{method} from above throws the forbidden Exception\index{exception} \lstinline|java.lang.RuntimeException|\index{java.lang!RuntimeException}. We want to fix that and we search the documentation for already existing Exceptions\index{exception} that fits. Among others, we find the class \lstinline|javax.xml.crypto.dsig.TransformException|\index{javax.xml.crypto.dsig!TransformException} \autocite{ORACLE_DOC_TRANSFORMEXCEPTION_CLASS}; while the name \bdq{\texttt{TransformException}} sounds good, but our code is in no way related to XML processing or even some cryptology stuff so we should not use that Exception to signal the failure of our \lstinline|textToNumber()| method. The Exception of choice would be \lstinline|java.lang.NumberFormatException|\index{java.lang!NumberFormatException} \autocite{ORACLE_DOC_NUMBERFORMATEXCEPTION_CLASS} that is also thrown by the \lstinline|Integer::parseInt| method.






The common recommendation is to create your own, domain-specific Exceptions instead of re-using already existing Exception classes – see below. Actually, the only Exceptions excluded from this are java.lang.IllegalArgumentException, java.lang.IllegalStateException, and perhaps java.lang.UnsupportedOperationException as well.
An IllegalArgumentException (and Exceptions that are derived from it) is usually thrown when an argument check fails (argument is null when it is not allowed to be null, a divisor is zero, an age is negative, …), and it makes sense to always use it – although a domain may derive its own version from IllegalArgumentException.
Domain Specific Exceptions
Each domain (component, module, area of operation, name it …) should have its “Domain specific Exception” – or, more precisely, its own, specific hierarchy of unchecked or checked (or both) Exceptions. Internally, the domain could have also another hierarchy of checked Exceptions, but these should not get visible outside that domain.
Except IllegalArgumentException and exceptions derived from that (“validation exceptions”), a domain or component may never emit an Exception that is not defined by itself (except NPE, but this may not occur anyway). All other Exceptions (Exceptions, not Errors!) must be caught and handled or wrapped into an Exception specific to the current domain. The validation exceptions are thrown only by the API methods of the domain as response of an invalid argument to that method. An IllegalArgumentException that is thrown by another component used by this domain must be wrapped and delegated or handled internally.
To illustrate that, we assume that we defined the domain (or component) “Printing” (in package com.pega) that is responsible for the generation of documents for paper-like output and their distribution. It defines the exception class PrintingException (or PrintingRuntimeException, if you like that better), derived from java.lang.RuntimeException. It will directly extend RuntimeException here, or if “Printing” would be part of another domain (if it has a parent component), it will extend that component’s domain specific exception. Perhaps you have defined an application exception (like our PRRuntimeException), so the domain specific exceptions will extend that.
This “domain root exception” PrintingException allows all consumers of the respective component to catch all exceptions from that component with a single catch:
try
{
    usePrinting();
}
catch( final PrintingException e )
{
    final var message = """
       printLetter(): Printing Operation failed!
       Recipient Id: %s
       Origin: %s
       """.formatted( recipientID, dataSource );
    oLog.error( message, e );
    terminate();
}
Alternatively, you wrap message and Exception into the appropriate domain specific Exception of this consumer and throw that:
try
{
    usePrinting();
}
catch( final PrintingException e )
{
    final var message = """
       Printing Operation failed!
       Recipient Id: %s
       Origin: %s
       """.formatted( recipientID, dataSource );
    throw new <MyDomainSpecificException>( message, e );
}
In this sample, recipientID and dataSource enhance the information that is already provided with the PrintingException object; probably the “Printing” component was called with the data itself (not with the dataSource), and the recipient is not part of the generated document.
The “Printing” domain now defines problem specific exceptions derived from PrintingException.

This allows the consumers to respond to specific problems in a specific manner: if the “Printing” component throws a PrintingDeviceException, the consumer knows that there is a problem with the output device and can send an email to the support staff to inform them that the printer needs maintenance. A PrintingDataException tells the consumer that the data it provided was faulty, and a PrintingTooMuchDataException (derived from PrintingDataException) signals that too much data should be printed. An InvalidPrintingFormatException would indicate that the output description is not valid, and a NoOutputDeviceException tells that there is no printer with the given name.
When the consumer wants to respond to all these exceptions differently (and it should do so! At least it should provide different data to the logging), the code snippet above needs to be changed to look like that below. Of course, we could always provide all information to an instance of PrintingException and leave it to the support engineer to identify what is relevant. But that would also mean that any kind of self-healing would get unnecessary complex.

try
{
    usePrinting();
}
catch( final NoOutputDeviceException e ) // Self-Healing
{
    final var message = """
       printLetter(): Printing Operation failed because device does not exist
       Attempting again with alternative device
       Recipient Id: %s
       Origin: %s
       Missing Device: %s
       Alternative Device: %s
       """.formatted( recipientID, dataSource, missingDevice, getAlternativeDevice() );
    oLog.warn( message, e );
    retryPrinting();
}
catch( final InvalidPrintingFormatException e )
{
    final var message = """
       Printing Operation failed due to an invalid format specification
       Specification: %s
       """.formatted( format );
    throw new IllegalArgumentException( message, e );
}
catch( final PrintingTooMuchDataException e )
{
    final var message = """
       Printing Operation failed because of too much data
       Recipient Id: %s
       Origin: %s
       Data Size: %d records
       """.formatted( recipientID, dataSource, dataSize );
    throw new <MyDomainSpecificException>( message, e );
}
catch( final PrintingDataException e )
{
    final var message = """
       Printing Operation failed because of invalid data
       Recipient Id: %s
       Origin: %s
       """.formatted( recipientID, dataSource );
    throw new <MyDomainSpecificOutputProblemException>( message, e );
}
catch( final PrintingException e )
{
    final var message = """
       Printing Operation failed with an unspecified Exception
       Recipient Id: %s
       Origin: %s
       """.formatted( recipientID, dataSource );
    oLog.error( message, e );
    abort();
}
The “Printing” component signals that it cannot proceed by throwing a PrintingException or one of the exceptions derived from it. What happens as response to that exception is on the consumer for “Printing”: it could terminate itself, too, it could attempt to print again, it can choose an alternative mean to get the documents out, or it could even ignore the problem … or it delegates the whole thing to its own consumers, by just wrapping the caught PrintingException into its own domain specific exception:
This decision (“Either Log (and Terminate) or (Wrap and) Throw”) must be made for each and every call.
Using an unchecked hierarchy (based on java.lang.RuntimeException) means that the consumer can completely ignore the exception handling for the PrintingException (and leave it to its own consumers). It also means that the API methods for the PrintingComponent do not need to have a throws clause, to declare the thrown Exceptions. A sample for this approach can be found with the java.util.regex package (the regular expression handling in Java). The other approach, to use checked exceptions (java.lang.Exception), can be found in java.io that defines IOException as the “domain specific exception” and derives exceptions like FileNotFoundException from that. Checked exceptions enforce that each call is somehow wrapped in a try-catch block, otherwise the code does not compile. This forces attention for the error handling on each caller but can bloat the code. One can say bloated code is much better than code without proper error handling, but when it is too bloated, the code is difficult to read and prone to get broken when amended.
The parameter lists of the constructors for the domain specific Exceptions are not limited to the regular Throwable parameter for the cause and the String parameter for the message. If there is parameter that is required by each API method of the domain, this could be also a mandatory parameter for the domain specific Exception. For the “Printing” domain/component this could be the id for the output device. In that case, that id is no longer required to be part of the Exception message.
Some words regarding the calls to terminate() vs. abort()in the sample code above: assume we use the code snippet in a loop, where each cycle uses different data to be printed. A PrintingDataException means that the current data is somehow corrupted, but we can proceed with the loop (terminate the current process and load the next dataset), but an InvalidPrintingFormatException or a PrintingDeviceException means that no printing is possible at all and we abort the loop by throwing the wrapper exception, and for the unspecific PrintingException, the consumer aborts.
There are libraries that use “Error” (domain specific implementations of java.lang.Error) to signal the abort condition, while “Exception” means, “you can try again with different (valid) input data”. It is a widely accepted best practice not to catch Error, so we should not adopt this approach.
Throwing Exceptions
That throwing an Exception will have its price and that it should be done with some care was discussed already above. Although not explicitly said in the chapter above, it should be obvious that the context for the Exception must be provided in the Exception message, similar as for the log message. Only the method’s name is not needed as this is part of the stack trace stored with the Exception itself.
Not all Exceptions that are emitted by a method are explicitly thrown by it! A RuntimeException (an unchecked Exception) thrown by a method called by the current method may bubble up to the caller if that current method, same as a checked exception that is declared to be thrown by the current method.
To allow a specific handling of the emitted Exceptions, each method must declare all of them (including the unchecked ones!), in a detailed list; a mere “throws Exception” is wrong!
Your method should fail fast, if at all. This means that all checks should be placed to the beginning of a method, and that an operation is aborted immediately after an unrecoverable error condition is identified.
This also means that each method (in Pega, this includes the Rules) must check immediately its arguments whether they are valid, and that it must throw a “validation exception” (an Exception that is derived from IllegalArgumentException) in case that check fails.
Catching Exceptions
First, you never catch java.lang.Throwable, java.lang.Exception, java.lang.RuntimeException and java.lang.Error.
There must be a dedicated catch block for each Exception that you want to handle; if the catch block does the same for multiple Exceptions, you list them in the catch clause:
catch( final ABCException | DEFException | XYZException e ) { … }
If you write an Exception to the log with a level of FATAL or ERROR, you must terminate or abort the current operation; by no means, you throw after you wrote a log message with log level ERROR or FATAL!
Catching an exception, logging it with a lower log level (WARN, INFO, DEBUG or TRACE), and throwing it again (wrapped or directly) can make sense, but you should think about it carefully.
But sometimes it makes sense to catch an Exception that can neither be handled nor logged: you should do it when you can enhance the information that comes with the Exception with the current context. You must do it when the exception is leaving the current domain (remember, a domain should only emit its domain specific Exceptions).
As a rule of thumb, each method that has parameters should catch all Exceptions and enhance them with the arguments it was called with, as shown above. Usually, you will do this by wrapping the caught exception, but you can also consider a design for your domain specific Exception that allows you to store additional information to the existing Exception instance.2
Wrapping looks like this:
try
{
    …
}
catch( final DomainException e ) 
{
    throw new DomainException( """
       Printing Operation failed
       Recipient Id: %s
       Origin: %s
       """.formatted( recipientID, dataSource ), e;
}
If the domain specific exception would allow to add information to it, that could look like this:
try
{
    …
}
catch( final DomainException e ) 
{
    e.addInfo( "Message", "Printing Operation failed" );
    e.addInfo( "Recipient Id", recipientID );
    e.addInfo( "Origin", dataSource ); 
    throw e;
}
“Handling an Exception” usually means that the error condition that caused the Exception is removed or healed. Whether or not the Exception is logged depends on the circumstances, but it must be obvious that it will never be logged with ERROR or FATAL when the error condition has gone!
Exceptions may never be swallowed. This means that code like this is inherently wrong or illegal:
\begin{lstlisting}
try
{
    …
}
catch( final AnException _ ) {}
\end{lstlisting}
Nonetheless, it is possible that you do not need to handle a caught Exception, either because the system is not in an error state despite the Exception (perhaps because you already provided a default value that is used in case the retrieval of a config value fails), or the Exception is declared for a method called from the try block, but cannot be thrown by that method. You must provide at least a comment in the otherwise empty block:
\begin{lstlisting}
catch( final AnException _ ) { /* Deliberately ignored */ }
\end{lstlisting}
Or, better, with an explanation: 
\begin{lstlisting}
catch( final AnException _ ) { /* Default value is used in case of failure */ }
\end{lstlisting}
In the second case, you should consider throwing a dedicated Exception:
catch( final AnException e ) { throw new ImpossibleExceptionError( e ); }
This protects the code a little bit from changes to the implementation of the method that previously could not throw the exception. At least the Exception will not go unnoticed when thrown.
Another recommendation regarding Error handling is that it must not obscure the essentials. When you write the code that checks for an error condition or handles an error condition, come to the point fast, otherwise the readability of the code suffers.
The finally Block
The complete syntax for the Exception handling is try-catch-finally:
\begin{lstlisting}
try
{
    …
}
catch( final AnException e ) { … }
finally
{
    // Some code
    …
}
\end{lstlisting}
The JVM guarantees that the code in the finally block will be executed under all circumstances, no matter how the try block is left.3
This is used for code like this:
theLock.lock(); // Retrieve the lock
try
{
    // Do something that needs synchronisation
    …
}
catch( final AnException e ) { … }
finally
{
    theLock.unlock(); // Release the lock
}
Unfortunately, we had INCs regarding locks that were held by already terminated threads. The interpretation was that there must be a bug in the JVM that, under some circumstances, allows a thread to terminate without executing a finally block.
But the solution was much more profane: there was code before the call to Lock::unlock that threw an exception so that the finally block was left before unlock() was called at all. 
The specification says that the code in the finally block is always executed, but it does not state that the code block is always left on its end or that the code is executed completely.
Consequently, all the housekeeping chores should be done at the very beginning of the finally block, with those calls first that cannot fail with an exception.
Exception Handling in Thread
The most often heard excuse for a catch phrase for java.lang.RuntimeException, java.lang.Exception or even java.lang.Throwable is that Exceptions thrown while executing a thread that is not the main thread may not go unnoticed. So you see code like this:
public final void run()
{
    try
    {
        // Do whatever this thread must do …
    }
    catch( final Exception e )
    {
        oLog.error( "No matter what went wrong, but I will log the exception", e );
    }
} 
Unfortunately, this absolves the programmer of the duty to properly handle any Exceptions thrown within the try block; the catch block at the end will take care of it.
The best practice is to set the UncaughtExceptionHandler for the thread, either through the ThreadFactory, through the ThreadGroup or individually. If your thread requires a special handling for an uncaught Exception, this looks like below; the sample assumes that the thread comes from a pool and is reused:
public final void run()
{
    final var currentThread = Thread.currentThread();
    final var uncaughtExceptionHandler = currentThread.getUncaughtExceptionHandler();
    currentThread.setUncaughtExceptionHandler( (t,e) ->
    {
        /*
         * The special exception treatment goes here …
         *

        // Finally log the exception
        oLog.fatal( """
            Thread "%s" was aborted with an exception!\
            """.formatted( t.getName() ), e );
    } );
 
    try
    {
        // Do whatever this thread must do …
    }
    finally
    {
        currentThread.setUncaughtExceptionHandler( uncaughtExceptionHandler );
    }
}
But usually a thread should use the UncaughtExceptionHandler that is assigned through the factory or group, as already said, or you even define a default handler. All this needs to be regulated by a proper guide to Exception Handling that needs to be defined.























\subsection{General Exception Handling}\label{sec:GeneralExceptionHandling}





















A method is only allowed to handle an exception if it is able to restore the program state, otherwise it has to delegate the exception. Delegation can be done either by ignoring the exception – in this case it will be propagated to the caller automatically (it “bubbles up”) – or by catching and wrapping it, as described in chapter \tqfullref{sec:ThrowingExceptions}.

To handle an exception you first have to catch it.

Catching \lstinline|java.lang.Exception| is discouraged in most cases; it \textit{can} be acceptable when you need to wrap all checked exceptions into a runtime exception or into an error:
\begin{lstlisting}
public final void run()
{
    try
    {
        …
    }
    catch( final Exception e )
    {
        throw new WrapperException( e );
    }
}   // run()
\end{lstlisting}

You should absolutely avoid to catch \lstinline|java.lang.Error| and all of its subclasses, \lstinline|java.lang.RuntimeException|, and of course, you should not never ever catch \lstinline|java.lang.Throwable|\index{java.lang!Throwable}!\footnote{But no rule without an exception: refer to chapter \tqvref{sec:ExceptionHandlingInThreads}.}

Provide a \lstinline|catch| block for each and every exception that you want to handle:
\begin{lstlisting}
public final void run()
{
    try
    {
        …
    }
    catch( final FileNotFoundException e )
    {
        throw new ApplicationError( "Cannot find file", e );
    }
    catch( final IOException e )
    {
        throw new ApplicationError( "An error occurred on processing the input file", e );
    }
    catch( final NumberFormatException e )
    {
        throw new ApplicationError( "Input Data is corrupted", e );
    }
    catch( final PatternSyntaxException e )
    {
        throw new ApplicationError( "Pattern invalid", e );
    }
}   // run()
\end{lstlisting}

Sometimes when a large number of exceptions has to be handled with all the same code, you may be tempted to just catch \lstinline|java.lang.Exception|. As already said: don't do it! This approach will definitely cause problems as it may handle runtime exceptions in the same \lstinline|catch|-block as the checked exceptions, and usually this is not intended.

A simplificatin for the sample above could look like this:
\begin{lstlisting}
public final void run()
{
    try
    {
        …
    }
    catch( final IOException | NumberFormatException | PatternSyntaxException e )
    {
        throw new ApplicationError( "Processing failed", e );
    }
}   // run()
\end{lstlisting}
Still only the declared exceptions will be handled in the \lstinline|catch|-block, but unexpected exceptions will now bubble up\footnote{The exception class \lstinline|java.io.FileNotFoundException| is not listed because its superclass \lstinline|java.io.IOException| is already on the list.}.

Sometimes an exception indicates a regular program state (although this would indicate poor design in most cases)\footnote{Nevertheless, you can find this even in the Java Runtime Library~… but no one has ever said that the authors of the JDK had never made doubtful decisions.}. Usually you will avoid this situation by e.g. checking the operands before an operation, but a very popular example where this is not practicable is the conversion of a string into a numerical value:
\begin{lstlisting}
…
var amount = 0.0;
final var requestParam = request.getParameter( "amount" );
try
{
    if( nonNull( requestParam ) )
    {
        amount = Double.parseDouble( requestParam );
    }
}
catch( final NumberFormatException ignored )
{
    //---* If amount in request is not numeric, set it to 0.0 *------
    amount = 0.0;
}
…
\end{lstlisting}

As shown in the example, a comment is required in the catch clause, because you do not further propagate the exception neither do you log it.

It is also possible that you checked the preconditions, but you still have to deal with some checked exceptions:
\begin{lstlisting}
…
final var file = new File( … );
if( file.exists() )
{
    try( final var input = new FileInputStream( file ) )
    {
        …
    }
    catch( final FileNotFoundException e )
    {
        throw new UnexpectedExceptionError( "File was there when checked!", e );
    }
}
…
\end{lstlisting}

A sample for the class \lstinline|UnexpectedExceptionError| can be found in \autocite{TQUADRAT_ORG_FOUNDATION_UNEXPECTEDEXCEPTIONERROR}.

Some APIs in the JDK declares a checked exception, but that will never be thrown for the given arguments by any means. Most famous for this are the conversions from byte arrays to Strings and back:
\begin{lstlisting}
…
final String s;
try
{
    s = new String( byteArray, "UTF-8" );
}
catch( final UnsupportedEncodingException e )
{
    throw new ImpossibleExceptionError( "UTF-8 must exist!", e );
}
…
\end{lstlisting}

According to \autocite{ORACLE_DOC_CHARSET:StandardCharsets}, the UTF-8 encoding is mandatory for each JDK/JVM implementation, therefore the code sequence above will \textit{never} throw an \lstinline|UnsupportedEncodingException|. Nevertheless, it is not allowed to swallow the exception; instead an error should be thrown in the \lstinline|catch| clause, indicating the error condition. In the improbable case UTF-8 is not supported by the current Java runtime environment, the problem will surface, and the program will not die silently without any hint about the reason.\footnote{If, of course, there is no catch for \lstinline|Throwable| somewhere on the path.} It also protects you from malicious changes to the name of the character encoding: a none existing \lstinline|Charset| will cause noise.

Of course, the problem could be solved differently, by using an instance of \lstinline|java.nio.charset.Charset|\autocite{ORACLE_DOC_CHARSET_CLASS} as the argument to the \lstinline|String| constructor, instead of the name of the character encoding. For this case, the constructor does not declare an \lstinline|UnsupportedEncodingException|. 

Another example is the interface \lstinline|java.lang.Appendable|\autocite{ORACLE_DOC_APPENDABLE_INTERFACE}: it is implemented by \lstinline|java.io.Writer| and \lstinline|java.io.PrintStream|, but also by \lstinline|java.lang.StringBuilder| (and a few more). The method \lstinline|java.lang.Appendable::append| declares to throw an \lstinline|IOException|, but the implementation of that method for \lstinline|StringBuilder| will never throw it and does not even declare it.
\begin{lstlisting}
…
final Appendable appendable = new StringBuilder();
try
{
    appendable.append( "Some Text" );
}
catch( final IOException e )
{
    throw new ImpossibleExceptionError( "Append to StringBuilder", e );
}
…
\end{lstlisting}

Again, a sample for \lstinline|ImpossibleExceptionError| can be found in \autocite{TQUADRAT_ORG_FOUNDATION_IMPOSSIBLEEXCEPTIONERROR}.

Sometimes a method wants to act upon an error condition without handling the exception. In this case, it may catch the exception, does whatever necessary, then throw it again. An example:
\begin{lstlisting}
public final void myMethod() throw IOException
{
    final Predicate<String> isNetworkError = m_NetworkErrorPattern.asMatchPredicate();
    …
    try
    {
        …
    }
    catch( final IOException e )
    {
        final var message = e.getMessage();
        if( nonNull( message ) && isNetworkError.test( message ) ) ++networkErrors;
        throw e;
    }
    …
}   // myMethod()
\end{lstlisting}

It is usually not a good idea to log an exception and to rethrow or to wrap it afterwards, because this might cause the exception to be logged twice (or even more often). More details about logging can be found in chapter \tqvref{sec:Logging}.

Finally, you handle the exception somehow, by assigning a default value\footnote{see the sample about the conversion from string to number, above.}, or by performing an alternative action. In any case, you need to provide a comment in the \lstinline|catch| why you do not further delegate the handling of the exception.

If an exception was not caught in your code, it ”bubbles up” until it reaches the main method for the current thread (that is \lstinline|main()| for the main or program thread, and \lstinline|run()|\autocite{ORACLE_DOC_RUNNABLE:run} for any other thread). If this method will not handle the exception, the thread will be aborted. If this thread is the last non-deamon thread, the JVM will be terminated, too.

%==============================================================================
% End of file!!
%==============================================================================

